\documentclass[preview]{standalone}

\usepackage{amsmath}
\usepackage{amssymb}
\usepackage{stellar}
\usepackage{definitions}

\begin{document}

\id{primary-modules}
\genpage

\section{Primary elements and components}

\begin{snippetdefinition}{primary-module-definition}{Primary elements, modules, and components}
    Let \(A\) be a \principalidealdomain, let \(M\) be an
    \(A\)-\module, and let \(p\in A\) be a \primeelement.
    \begin{enumerate}
        \item An element \(v\in M\) is \emph{\(p\)-primary} if
        \[
            p^n v=0_M
        \]
        for some \(n\in\naturalnumbers\). Equivalently,
        \(p^n\in\moduleann(v)\) for some \(n\).
        \item The \module \(M\) is \emph{\(p\)-primary} if every one of
        its elements is \(p\)-primary. It is \emph{primary} if it is
        \(p\)-primary for some \primeelement \(p\).
        \item The \emph{\(p\)-primary component} of \(M\) is
        \[
            M_p
            =
            \{v\in M\suchthat p^n v=0_M
            \text{ for some }n\in\naturalnumbers\}.
        \]
    \end{enumerate}
\end{snippetdefinition}

\begin{snippetproposition}{primary-component-submodule}{The primary component is a submodule}
    Let \(M\) be an \(A\)-\module. Then
    \[
        M_p\submoduleleq\moduletor(M).
    \]
    Moreover, \(M_p\) is the largest \(p\)-\primarymodule contained in
    \(M\).
\end{snippetproposition}

\begin{snippetproof}{primary-component-submodule-proof}{primary-component-submodule}{The primary component is a submodule}
    Since \(p^0 0_M=0_M\), one has \(0_M\in M_p\). If \(v,w\in M_p\),
    choose \(n,m\in\naturalnumbers\) such that \(p^nv=0_M\) and
    \(p^mw=0_M\). With \(r=\max\{n,m\}\),
    \[
        p^r(v-w)=0_M.
    \]
    If \(a\in A\), then \(p^n(av)=a(p^nv)=0_M\). Hence
    \(M_p\submoduleleq M\).

    Every nonzero \(v\in M_p\) is annihilated by some nonzero power of
    \(p\), while \(0_M\) is torsion because \(\moduleann(0_M)=A\).
    Thus \(M_p\subseteq\moduletor(M)\). Finally, every \(p\)-primary
    \submodule of \(M\) consists entirely of elements of \(M_p\), which
    proves maximality.
\end{snippetproof}

\section{The socle}

\begin{snippetdefinition}{primary-module-socle-definition}{Socle of a primary module}
    Let \(M\) be a \(p\)-\primarymodule. Its \emph{socle} is
    \[
        \modulesocle_p(M)
        =
        \{v\in M\suchthat pv=0_M\}.
    \]
    More generally, for any \(A\)-\module \(M\), the same formula defines
    the \(p\)-socle of its \primarycomponent \(M_p\).
\end{snippetdefinition}

\begin{snippetproposition}{primary-module-socle-vector-space}{The socle as a vector space}
    Let \(K=A/(p)\). Then
    \[
        \modulesocle_p(M)\submoduleleq M_p,
    \]
    and \(\modulesocle_p(M)\) is naturally a \(K\)-\vectorspace under
    \[
        (a+(p))v=av.
    \]
\end{snippetproposition}

\begin{snippetproof}{primary-module-socle-vector-space-proof}{primary-module-socle-vector-space}{The socle as a vector space}
    The zero element lies in the socle. If \(v,w\) lie in it and
    \(a\in A\), then
    \[
        p(v-w)=pv-pw=0_M,
        \qquad
        p(av)=a(pv)=0_M.
    \]
    Hence it is a \submodule.

    Since \(A\) is a PID, the \ideal \((p)\) is maximal and \(K\) is a
    \field. To check that the displayed scalar action is well-defined,
    suppose \(a'=a+bp\). For \(v\in\modulesocle_p(M)\),
    \[
        a'v=(a+bp)v=av+b(pv)=av.
    \]
    The vector-space axioms now follow from the module axioms.
\end{snippetproof}

\end{document}
